%================================
% Vocabulary & Language Focus — XB3 (Unit 3)
% 编译方式: xelatex main.tex (需运行两遍以生成目录)
%================================
\documentclass[../main.tex]{subfiles}

\begin{document}

\section{Unit 3}

\vw{breadth}{/bredθ/}{n.}

\sense{n.}{宽度}
\col{the breadth of sth}
\ex{The breadth of the river at this point is nearly a kilometre.}

\sense{n.}{广泛}
\col{be breathtaking breadth and depth 广度与深度}
\ex{The course covers an impressive breadth and depth of modern European history.}
\col{the breadth of knowledge/experience}
\ex{His breadth of knowledge is truly remarkable for someone so young.}

\begin{gram}{breadth / width / depth / length 的"维度"细分}
\textbf{length}（长）；\textbf{width / breadth}（宽——前者通用，后者偏文学/抽象）；\textbf{depth}（深）；\textbf{height}（高）。breadth 多与 depth 联用，构成"广度与深度"对偶：\textit{the breadth and depth of his learning}。
\end{gram}

\begin{origin}
古英语 \textit{brǣdu}（宽度、广度），是 \textit{brād}（宽的）的抽象名词。同源词 \textbf{broad}（宽的）、\textbf{broaden}（使变宽）、\textbf{abreast}（并排；并肩；\textit{a-} 在 + \textit{breast} 胸 $\to$ 胸对胸$\to$肩并肩）。breadth 与 broad 同源——形容词名词化，\textit{breath}（呼吸）来自不同词根 \textit{bræs} 气息。
\end{origin}

\vw{convince}{/kənˈvɪns/}{vt.}

\sense{vt.}{使确信，使相信}
\col{convince sb of sth}
\ex{He managed to convince the jury of his innocence.}
\col{convince sb that ...}
\ex{Scientists have convinced the public that climate change is real.}

\sense{vt.}{说服，劝说}
\col{convince sb to do sth}
\ex{I couldn't convince her to come with us.}

\begin{gram}{convince / persuade 的"说服"细分}
\textbf{convince}（使相信——侧重结果：\textit{convince sb that...} 使某人相信……）；\textbf{persuade}（说服——侧重过程：\textit{persuade sb to do sth}）。考试常考：\ding{55} \textit{convince sb to do sth}；\checkmark\ \textit{convince sb that...} 或 \textit{persuade sb to do sth}。但现代美式英语中 convince sb to do sth 已普及。
\end{gram}

\begin{origin}
16 世纪借自拉丁语 \textit{convincere}（证明有罪；使确信；反驳）：\textit{con-}（完全）+ \textit{vincere}（征服）。同源词 \textbf{convincing}（令人信服的）、\textbf{conviction}（确信；定罪）、\textbf{convince}、\textbf{vanquish}（征服；\textit{vincere} 的古法语形式）、\textbf{victor}（胜利者）、\textbf{victory}（胜利）、\textbf{evince}（表明；\textit{e-} 出 + vincere $\to$ 战胜到出$\to$ 证明）、\textbf{province}（省份；\textit{pro-} 前 + vincere $\to$ 前方征服之地）。
\end{origin}

\vw{voyage}{/ˈvɔɪɪdʒ/}{n. \& vi.}

\sense{n.}{航海，航行，航天}
\col{go on/make a voyage}
\ex{Columbus made his famous voyage to the Americas in 1492.}
\col{a voyage of discovery}
\ex{The expedition was a voyage of discovery into the unknown interior of the continent.}

\sense{vi.}{航行，远行}
\col{voyage to/across sth}
\ex{The Mayflower voyaged across the Atlantic to the New World.}

\begin{lex}{voyage / journey / trip / travel / tour 的"旅行"细分}
\textbf{voyage}（长途海上/空中航行；\textit{a transatlantic voyage}）；\textbf{journey}（长途陆上/空中旅行；尤指单程：\textit{a three-day journey}）；\textbf{trip}（短途旅行；日常往返：\textit{a day trip}）；\textbf{travel}（旅行——泛指；作动词：\textit{travel abroad}）；\textbf{tour}（游览；周游：\textit{a round-the-world tour}）。voyage 最具航海/航天的史诗感。
\end{lex}

\begin{origin}
13 世纪借自古法语 \textit{voiage}，拉丁语 \textit{viaticum}（旅行用品；旅行费），源于 \textit{via}（路）。同源词 \textbf{viaduct}（高架桥；\textit{via} 路 + \textit{ducere} 引$\to$ 引路$\to$桥）、\textbf{deviate}（偏离；\textit{de-} 离 + via）、\textbf{obvious}（明显的；\textit{ob-} 在前 + via $\to$ 在路前$\to$挡在路前的）、\textbf{previous}（先前的；\textit{prae-} 在前 + via $\to$走在前面的）、\textbf{trivial}（琐碎的；\textit{tri-} 三 + via $\to$三岔路口的$\to$到处都见到的$\to$平庸的）。"航海"字面"上路"——古罗马人把旅行费叫 viaticum。
\end{origin}

\vw{fleet}{/fliːt/}{n.}

\sense{n.}{舰队}
\col{an aircraft/naval fleet}
\ex{The British fleet set sail for the South Atlantic at dawn.}
\col{fleet admiral 海军五星上将}
\ex{Fleet Admiral Chester Nimitz commanded the US Pacific Fleet during World War II.}

\sense{n.}{捕鱼船队；车队}
\col{a fishing fleet}
\ex{The harbour is home to a large fishing fleet.}
\col{a fleet of taxis/buses/cars}
\ex{A fleet of taxis was waiting outside the station.}

\begin{origin}
古英语 \textit{flēot}（船；河流中的浅滩），与 \textit{flēotan}（漂浮、流）同源。同源词 \textbf{fleet}（船队；快速移动；\textit{a fleet-footed runner} 快脚的）、\textbf{float}（漂浮）、\textbf{flood}（洪水）、\textbf{flotilla}（小舰队；西班牙语 \textit{flota} 船队）。fleet 在古英语中指"船"或"浅滩"，现代指"船队"是延续本义。
\end{origin}

\vw{navigate}{/ˈnævɪɡeɪt/}{vt. \& vi.}

\sense{vt.}{航行；导航}
\col{navigate a ship/plane}
\ex{The captain navigated the ship safely through the narrow channel.}
\col{navigate (by sth)}
\ex{Early sailors navigated by the stars.}

\sense{vi.}{导航}
\col{navigate through/across sth}
\ex{We had to navigate through heavy traffic to get to the airport on time.}

\begin{gram}{navigate the internet / a website：抽象"导航"用法}
\textbf{navigate the web/internet}（在网上浏览）、\textbf{navigate a complex situation}（在复杂局势中摸索前进）、\textbf{navigate a maze}（走迷宫）。navigate 的抽象用法皆源自物理航行——在未知领域中"找方向"。
\end{gram}

\begin{origin}
16 世纪借自拉丁语 \textit{navigatus}（航行的），是 \textit{navigare}（航行、驾驶）的过去分词：\textit{navis}（船）+ \textit{agere}（引导、行动）。同源词 \textbf{navigation}（航行；导航）、\textbf{navigator}（航海者；领航员）、\textbf{naval}（海军的；船的）、\textbf{circumnavigate}（环航；\textit{circum-} 周围 + navigare）、\textbf{navy}（海军）、\textbf{nausea}（晕船；\textit{nausia} 船病）。"航行"字面"使船动"——agere 引导与 navis 船的合成。
\end{origin}

\vp{on behalf of}

\sense{短语}{代表（或代替）；为帮助；为了}
\col{on behalf of sb}
\ex{I'm writing on behalf of my mother, who is unwell and cannot attend the meeting.}
\ex{The lawyer spoke on behalf of her client at the hearing.}
\col{on sb's behalf = on behalf of sb}
\ex{Please donate generously on behalf of the victims of the earthquake.}

\begin{gram}{on behalf of / in behalf of：英美用法差异}
英式英语严格区分：\textbf{on behalf of sb}（代表某人；代替某人——代表身份）；\textbf{in behalf of sb}（为了某人；为某人利益——利益立场）。美式英语中两者常混用，多用 on behalf of。behalf 本身无独立含义，是 \textit{be-} + \textit{half}（"在旁边"）的合成。
\end{gram}

\vw{load}{/ləʊd/}{vt. \& vi. \& n.}

\sense{vt. \& vi.}{装上，装入}
\col{load sth (into/onto sth)}
\ex{They loaded the luggage into the back of the car.}
\col{load a gun/camera/computer program}
\ex{She loaded the camera with a new roll of film.}

\sense{n.}{负载；承载量}
\col{a heavy load}
\ex{The truck was carrying a heavy load of timber.}
\col{a load of ... 一大堆}
\ex{We've got a load of work to finish before the weekend.}

\sense{n.}{大量}
\col{loads of sth}
\ex{There were loads of people at the concert.}

\begin{gram}{load / cargo / freight / burden 的"载"细分}
\textbf{load}（载；通用：\textit{a load of hay}）；\textbf{cargo}（货物；多指船/飞机所载：\textit{a ship full of cargo}）；\textbf{freight}（货运；多指铁路/公路运输：\textit{freight train} 货运列车）；\textbf{burden}（负担；常含情感/道德重负：\textit{a heavy burden of guilt}）。burden 的"重"是隐喻，load 的"重"是物理。
\end{gram}

\begin{origin}
古英语 \textit{lād}（路、行程；载货于路上），与 \textit{lith}（船的吃水线）相关。同源词 \textbf{lade}（装载；\textit{lade} + \textit{-en} $\to$ laden 装满的——\textit{a laden ship} 满载的船）、\textbf{laden}（装满的；\textit{heavily laden} 装载沉重的）、\textbf{unload}（卸货）。load 与 \textit{lode}（矿脉；\textit{lode-star} 指示星）形近但来源不同。
\end{origin}

\vw{compass}{/ˈkʌmpəs/}{n.}

\sense{n.}{罗盘，指南针}
\col{a magnetic compass}
\ex{Sailors have used the magnetic compass for navigation since the 12th century.}
\col{compass points 罗盘方位}
\ex{The four main compass points are north, south, east, and west.}

\sense{n.}{圆规}
\col{a pair of compasses}
\ex{The architect drew a perfect circle with a pair of compasses.}

\begin{gram}{compass / compasses 的可数性}
\textbf{compass}（罗盘）单数；\textbf{compasses}（圆规）复数。两种含义各有可数规则：\textit{a compass} 一只罗盘、\textit{a pair of compasses} 一副圆规。同形异义，词源各异。
\end{gram}

\begin{origin}
14 世纪借自古法语 \textit{compas}（圆规、量具），源于 \textit{compasser}（测量），拉丁语 \textit{compassare}（共同步测）：\textit{com-}（一起）+ \textit{passus}（步）。同源词 \textbf{compass}（罗盘——借步测距离的意象）、\textbf{encompass}（包含；包围；\textit{en-} 入 + compass $\to$ 走一圈$\to$包围）、\textbf{pace}（步；步速）、\textbf{pass}（经过；通行证）、\textbf{surpass}（超越）。"罗盘"字面"共同步测"——古罗马士兵用步测丈量土地。
\end{origin}

\vw{log}{/lɒɡ/}{n. \& vi. \& vt.}

\sense{n.}{航海日志，正式记录}
\col{keep a log}
\ex{The captain kept a detailed log of the ship's voyage.}
\col{a log book 日志簿}
\ex{Every car should have a complete log book of its service history.}

\sense{n.}{原木}
\col{a log fire 篝火}
\ex{They sat around a log fire and told stories late into the night.}
\col{sleep like a log 睡得死沉}
\ex{After the long hike, I slept like a log.}

\sense{vi. \& vt.}{伐木}
\col{log trees}
\ex{The workers logged the trees and floated them down the river.}
\col{logging 伐木业}
\ex{Logging has been a major industry in this part of the country for centuries.}

\begin{gram}{log in / log on / log out：电脑时代的新义}
\textbf{log in/on}（登录：\textit{log in to the system}）、\textbf{log out/off}（登出）。源自航海/航空日志（log book）——每一次操作都被"记录"下来。
\end{gram}

\begin{origin}
14 世纪中古英语 \textit{logge}（原木、大木块），可能与古挪威语 \textit{lag}（倒下的树）相关。同源词 \textbf{logger}（伐木工）、\textbf{logging}（伐木业）、\textbf{backlog}（积压；\textit{back} 后 + log $\to$ 后面堆着的木头）。"log"作为原木本是伐木业的产物，与"logarithm"（对数）形近但不同源。
\end{origin}

\vw{harmonious}{/hɑːˈməʊniəs/}{adj.}

\sense{adj.}{友好和睦的}
\col{a harmonious relationship/marriage}
\ex{They have worked hard to maintain a harmonious relationship with their neighbours.}

\sense{adj.}{协调的，和谐的}
\col{harmonious colours/sounds}
\ex{The garden is full of harmonious colours that complement each other beautifully.}

\begin{lex}{harmonious / amicable / cordial / friendly 的"和睦"细分}
\textbf{harmonious}（和谐的；强调整体协调感：\textit{a harmonious community}）；\textbf{amicable}（友善的；强调避免冲突：\textit{an amicable divorce} 友好离婚）；\textbf{cordial}（真诚热情的；\textit{a cordial welcome}）；\textbf{friendly}（友好的；通用）。harmonious 是最带"美学/音乐感"的"和"。
\end{lex}

\begin{origin}
15 世纪由 \textit{harmony}（和谐）+ \textit{-ous} 派生。同源词 \textbf{harmony}（和谐；和声）、\textbf{harmonize}（和谐；配和声）、\textbf{harmonica}（口琴）、\textbf{inharmonious}（不和谐的）。注意"和声"（music harmony）的"和"与"和谐"（peaceful harmony）的"和"在英语中皆 harmony。
\end{origin}

\vw{dismiss}{/dɪsˈmɪs/}{vt.}

\sense{vt.}{不予考虑，对……不屑一提}
\col{dismiss sth as sth}
\ex{He dismissed the report as nothing more than gossip.}
\col{dismiss a suggestion/idea}
\ex{The committee dismissed her suggestion without serious consideration.}

\sense{vt.}{去除，摒除}
\col{dismiss sth from one's mind}
\ex{Try to dismiss the thought from your mind and focus on the present.}

\sense{vt.}{解雇}
\col{dismiss sb for sth}
\ex{She was dismissed for repeated lateness.}
\col{be dismissed from one's job/post}
\ex{He was dismissed from his post after the scandal.}

\sense{vt.}{解散}
\col{dismiss the class/meeting/troops}
\ex{The teacher dismissed the class early because of the snow.}

\begin{gram}{dismiss / fire / sack / lay off 的"解雇"语体谱}
\textbf{dismiss}（解雇；正式：\textit{be dismissed from one's job}）；\textbf{fire}（解雇；口语美式：\textit{get fired}）；\textbf{sack}（解雇；口语英式：\textit{get the sack}）；\textbf{lay off}（暂时解雇/裁员：\textit{laid off workers}）。dismiss 较正式、原因常为失职；fire/sack 可无故；lay off 多指经济性裁员。
\end{gram}

\begin{origin}
15 世纪借自拉丁语 \textit{dismissus}（送走、遣散），是 \textit{dimittere}（送走、解散）的过去分词：\textit{dis-}（分开）+ \textit{mittere}（送）。同源词 \textbf{dismissal}（解雇；解散）、\textbf{dismissive}（轻蔑的；\textit{a dismissive wave of the hand}）、\textbf{missile}（导弹；"被发射之物"）、\textbf{mission}（任务；使命）、\textbf{emit}（发射；发出）、\textbf{submit}（提交；屈服）、\textbf{transmit}（传输；\textit{trans-} 越过 + mittere）、\textbf{promise}（承诺；\textit{pro-} 前 + mittere $\to$ 预先放出之言）、\textbf{admission}（承认；入场费）。
\end{origin}

\vw{proof}{/pruːf/}{n.}

\sense{n.}{证据，证明}
\col{proof of sth}
\ex{Can you provide proof of your identity?}
\col{proof that ...}
\ex{We have no proof that he was at the scene of the crime.}
\col{positive/conclusive proof}
\ex{The forensic evidence provided conclusive proof of the defendant's guilt.}

\begin{gram}{proof / evidence / testimony / witness 的"证"细分}
\textbf{proof}（证据；严格意义上"充分证明"：\textit{positive proof}）；\textbf{evidence}（证据；通用：\textit{circumstantial evidence} 间接证据）；\textbf{testimony}（证词；多指法庭/事件当事人陈述：\textit{give testimony}）；\textbf{witness}（证人；亲眼目击者：\textit{a witness to the accident}）。注意 evidence 不可数（统称"证据"）时可数（指"一份证据"）。
\end{gram}

\begin{origin}
13 世纪借自古法语 \textit{prove}（证据），古英语 \textit{prōf}，拉丁语 \textit{proba}（证明、证据），是 \textit{probare}（证明、检验）的派生形式。同源词 \textbf{prove}（证明；证实）、\textbf{proof}（证据；抗……的：\textit{watertight / fireproof}）、\textbf{approve}（赞同；\textit{ap-} 朝向 + probare $\to$ 证明为善$\to$ 赞同）、\textbf{improve}（改善；古义"证明为有利"——\textit{en-} 使 + \textit{prove} 旧形式）、\textbf{reprove}（责备）、\textbf{probable}（可能的；\textit{pro-} 试 + probare $\to$ 试过证明有效）。
\end{origin}

\vw{acknowledge}{/əkˈnɒlɪdʒ/}{vt.}

\sense{vt.}{承认（权威、地位）}
\col{acknowledge sb as sth}
\ex{She is widely acknowledged as one of the greatest writers of the century.}
\col{acknowledge sb's authority/status}
\ex{The new government refused to acknowledge the previous regime's treaties.}

\sense{vt.}{承认（属实）}
\col{acknowledge (that) ...}
\ex{He acknowledged that he had made a mistake.}
\col{acknowledge doing sth}
\ex{She acknowledged having been at the meeting.}

\sense{vt.}{(公开) 感谢}
\col{acknowledge sb's help/support}
\ex{I would like to acknowledge the help of all those who contributed to the project.}

\begin{gram}{acknowledge / admit / confess / recognize 的"承认"细分}
\textbf{acknowledge}（承认——公开/正式：\textit{acknowledge one's mistake} 公开认错）；\textbf{admit}（承认——不情愿的：\textit{admit one's guilt} 认罪）；\textbf{confess}（承认——含忏悔意味：\textit{confess to a crime} 忏悔罪行）；\textbf{recognize}（承认——认可/认出：\textit{recognize the new state} 承认新国家）。acknowledge 比 admit 更"体面"——承认事实但不必道歉。
\end{gram}

\begin{origin}
15 世纪由 \textit{ac-}（古法语 \textit{ac-} 形式 = ad-）+ \textit{knowledge}（知识；古英语 \textit{cnāwan} 知道 + \textit{-ledge} 抽象名词）合成。同源词 \textbf{acknowledge}、\textbf{knowledge}（知识；学问）、\textbf{acknowledgeable}（可承认的）、\textbf{gnaw}（咬；\textit{cn-} 知道 + \textit{gnaw} 咬——同源但义已分）。"承认"字面"知道"——知晓事实即承认事实。
\end{origin}

\vw{influential}{/ˌɪnfluˈenʃl/}{adj.}

\sense{adj.}{有很大影响力的，有支配力的}
\col{an influential politician/writer}
\ex{She is one of the most influential figures in modern art.}
\col{influential in sth}
\ex{Newton was highly influential in shaping our understanding of the physical world.}

\begin{lex}{influential / powerful / dominant / prestigious 的"有力"细分}
\textbf{influential}（有影响力的——常含褒义"靠思想/能力影响他人"：\textit{an influential thinker}）；\textbf{powerful}（强大的——通用：\textit{a powerful engine}）；\textbf{dominant}（占主导的——\textit{a dominant species}）；\textbf{prestigious}（有声望的——\textit{a prestigious university}）。influential 强调"影响他人决策/思想"。
\end{lex}

\begin{origin}
16 世纪由 \textit{influence}（影响）+ \textit{-ial} 派生。influence 14 世纪借自古法语 \textit{influence}（流入、星力、感应），中古拉丁语 \textit{influentia}（流入），源于 \textit{influere}（流入）：\textit{in-}（入）+ \textit{fluere}（流）。同源词 \textbf{influence}（影响；影响力）、\textbf{influenza}（流感——"被（星力）流入"）、\textbf{fluent}（流利的；\textit{fluere} 流$\to$ 语言流畅的）、\textbf{fluid}（流体的；流动的）、\textbf{flush}（脸红；冲水；\textit{fluere} 流变体）、\textbf{flux}（流动；变迁）、\textbf{confluent}（汇合的）。"影响"原义"星体流入人体的力"——中世纪占星术概念。
\end{origin}

\vw{insist}{/ɪnˈsɪst/}{vi. \& vt.}

\sense{vi.}{坚决要求}
\col{insist on/upon sth/doing sth}
\ex{She insisted on paying for the meal.}
\col{insist on sb doing sth}
\ex{They insisted on him staying for dinner.}

\sense{vt.}{坚持}
\col{insist that ... (should) ...}
\ex{She insisted that he (should) be present at the meeting.}
\ex{I insist that you (should) take a break.}

\sense{vi.}{坚持说}
\col{insist on sth}
\ex{He insisted on his innocence throughout the trial.}

\begin{gram}{insist that 后接 should 型虚拟式（subjunctive）}
\textbf{insist}、\textbf{suggest}、\textbf{demand}、\textbf{recommend}、\textbf{require} 等动词后接 that 从句时，谓语动词应用\textbf{should + 动词原形}或\textbf{动词原形}（美式英语常省略 should）：\textit{She insisted that he \textbf{be} present.}（\ding{55} \textit{insisted that he \textbf{was} present}——在"坚持要求"义下）。这是英语"虚拟式"的标志——虽然现代英语的 -should 形式已不严格使用，但语法写作常考。
\end{gram}

\begin{origin}
16 世纪借自拉丁语 \textit{insistere}（坚持、站立不动）：\textit{in-}（在上）+ \textit{sistere}（站立）。同源词 \textbf{insistence}（坚持）、\textbf{insistent}（坚持的；急切的）、\textbf{consist}（组成；\textit{con-} 一起 + sistere $\to$ 站在一起$\to$由……组成）、\textbf{persist}（坚持；\textit{per-} 贯穿 + sistere $\to$ 始终站着不动）、\textbf{resist}（抵抗；\textit{re-} 反 + sistere）、\textbf{assist}（帮助；\textit{as-} 旁 + sistere $\to$ 站在旁边）、\textbf{state}（状态；陈述；\textit{sistere} 站立$\to$ 站着的样子$\to$陈述）。
\end{origin}

\vw{route}{/ruːt/}{n.}

\sense{n.}{路线，路途}
\col{a bus/shipping/cycle route}
\ex{We took the scenic route along the coast.}
\col{en route 在途中}
\ex{The delegation is en route to Beijing for the trade talks.}

\sense{n.}{途径，渠道}
\col{a route to sth}
\ex{Hard work is often the surest route to success.}
\col{the route to power/wealth}
\ex{She took an unusual route to power, starting as a research scientist.}

\begin{lex}{route / path / road / way / course 的"路"细分}
\textbf{route}（路线；规划好的：\textit{the shortest route}）；\textbf{path}（小路；散步小径：\textit{a garden path}）；\textbf{road}（公路；通用：\textit{address}）；\textbf{way}（道路；抽象方法：\textit{the right way}）；\textbf{course}（方向；进程：\textit{the course of history}）。route 偏"已规划"，way 偏"通用"，path 偏"自然/隐喻"。
\end{lex}

\begin{origin}
13 世纪借自古法语 \textit{route}，拉丁语 \textit{rupta}（开辟的路），是 \textit{rompere}（打破）的阴性过去分词。同源词 \textbf{route}、\textbf{routine}（常规；\textit{route} + \textit{-ine} $\to$ 走老路$\to$常规）、\textbf{rout}（溃败；\textit{rompere} $\to$ 军队溃散$\to$ 溃逃）、\textbf{abrupt}（突然的；\textit{ab-} 离 + \textit{rompere} $\to$ 突然断开的）、\textbf{corrupt}（腐败的；\textit{cor-} 一起 + \textit{rompere} $\to$ 一起打破规则$\to$ 腐败的）、\textbf{disrupt}（中断；扰乱）、\textbf{erupt}（爆发；\textit{e-} 出 + \textit{rompere} $\to$ 破出$\to$火山爆发）、\textbf{rupture}（破裂；断裂）。
\end{origin}

\vw{financial}{/faɪˈnænʃl/}{adj.}

\sense{adj.}{财政的，金融的，财务的}
\col{a financial crisis/adviser}
\ex{The country is in the middle of a serious financial crisis.}
\col{financial support/aid 财政支持}
\ex{She provided financial support to her younger brother during his college years.}

\begin{gram}{financial / fiscal / monetary / economic 的"财经"细分}
\textbf{financial}（财务的；通用：\textit{financial problems}）；\textbf{fiscal}（财政的；专指政府/公共财政：\textit{fiscal policy} 财政政策）；\textbf{monetary}（货币的；专指货币体系：\textit{monetary policy} 货币政策）；\textbf{economic}（经济的；泛指：\textit{economic growth}）。考试常考：fiscal policy 财政政策，monetary policy 货币政策。
\end{gram}

\begin{origin}
19 世纪由 \textit{finance}（财政；金融）+ \textit{-ial} 派生。finance 15 世纪借自古法语 \textit{finance}（终结、支付），拉丁语 \textit{finantia}（付款），源于 \textit{finire}（终结、结束）。同源词 \textbf{finance}（财政；金融）、\textbf{finances}（财力；财务状况）、\textbf{financial}（财政的）、\textbf{fine}（罚款；精细的；终结；\textit{finis} 结尾）、\textbf{finish}（完成；\textit{finire} 的英语变体）、\textbf{finite}（有限的；\textit{finis} 终结$\to$ 有终点的）、\textbf{infinite}（无限的）、\textbf{define}（定义；\textit{de-} 下 + \textit{finire} $\to$ 划定界限$\to$定义）、\textbf{refine}（精炼；\textit{re-} 再 + \textit{finire} $\to$ 再终结$\to$ 提纯）。
\end{origin}

\vw{leak}{/liːk/}{vi. \& vt. \& n.}

\sense{vi.}{漏，渗漏}
\col{leak from/through sth}
\ex{Water was leaking from the pipe under the kitchen sink.}
\col{the roof leaks 屋顶漏水}
\ex{The roof leaks badly whenever it rains.}

\sense{vi.}{渗入}
\col{leak in}
\ex{Rain had leaked in through a crack in the window frame.}

\sense{vt.}{漏，泄露}
\col{leak sth (to sb)}
\ex{Somebody leaked the news to the press.}

\sense{n.}{漏洞}
\col{a leak in sth}
\ex{There is a leak in the pipe that needs fixing immediately.}
\col{a gas/oil leak}
\ex{The explosion was caused by a gas leak in the basement.}

\begin{lex}{leak / drip / seep / ooze 的"漏/渗"细分}
\textbf{leak}（漏——容器有缺口：\textit{the pipe is leaking}）；\textbf{drip}（滴——一滴一滴落：\textit{water is dripping from the tap}）；\textbf{seep}（渗——缓慢地透过小孔：\textit{water seeped into the ground}）；\textbf{ooze}（渗出——粘稠液缓慢渗出：\textit{blood oozed from the wound}）。leak 偏"未预期的"，drip 偏"有节奏的"，seep 偏"自然渗透"，ooze 偏"粘稠"。
\end{lex}

\begin{origin}
13 世纪中古英语 \textit{leken}（漏水），可能与古挪威语 \textit{leka}（滴下、漏）相关。同源词 \textbf{leaky}（漏的；不严密的）、古英语 \textit{hlec}（泄漏）、荷兰语 \textit{lek}（漏洞）。leak 与 \textit{leakage}（渗漏）同根，与 \textit{lick}（舔）形近但实不同源。
\end{origin}

\vw{panic}{/ˈpænɪk/}{n. \& vi. \& vt.}

\sense{n.}{惊恐，恐慌}
\col{in a state of panic}
\ex{The crowd was in a state of panic as the fire spread through the building.}
\col{panic buying/selling 恐慌性购买/抛售}
\ex{There was widespread panic buying of face masks at the start of the pandemic.}

\sense{vi. \& vt.}{(使) 惊慌，(使) 惊慌失措}
\col{panic over/about sth}
\ex{Don't panic over the exam — you've prepared well.}
\col{panic sb into doing sth}
\ex{The sudden noise panicked the birds into flying away.}

\begin{gram}{panic / scare / frighten / terrify 的"惊"强度谱}
\textbf{panic}（恐慌——失措的群体性惊恐：\textit{The crowd panicked}）；\textbf{scare}（吓——轻度：\textit{scare sb stiff} 吓呆了）；\textbf{frighten}（惊——通用：\textit{frightened of spiders}）；\textbf{terrify}（恐吓——最强：\textit{terrified of heights}）。panic 多含"无法理性应对"之意。
\end{gram}

\begin{origin}
17 世纪借自希腊语 \textit{Panikos}（属于潘神的），源于 \textit{Pan}（潘神——希腊神话中半人半羊的山林之神）。潘神栖息山林，其突然出现令旅人骤然惊恐——panic 因此得名。同源词 \textbf{panic-stricken}（惊慌失措的）、\textbf{panicky}（容易恐慌的）。希腊语 \textit{pan}（全、所有）与英语 \textbf{pan-}（全……）同源——Pan 本是"全希腊"的山林守护者。
\end{origin}

\vw{regardless}{/rɪˈɡɑːdləs/}{adv.}

\sense{adv.}{不顾，不加理会}
\ex{She knew it was dangerous, but she went on regardless.}
\ex{Whatever the reason, he refused to help, regardless.}

\begin{gram}{regardless 与 regardless of 的搭配差异}
\textbf{regardless}（adv. 不顾——单独使用，多修饰整句：\textit{He went on regardless}）；\textbf{regardless of}（prep. 不顾——后接名词/动名词：\textit{regardless of the cost} 不顾成本；\textit{regardless of what you say} 不管你说什么）。前者是状语，后者是介词短语。\textit{irregardless} 是非标准用法，应避免。
\end{gram}

\begin{origin}
15 世纪由 \textit{regard}（注意；尊重）+ \textit{-less}（无）合成。同源词 \textbf{regard}（注意；尊重；\textit{give my regards to...} 代我向……问好）、\textbf{regardless}（不顾）、\textbf{disregard}（忽视；不顾）、\textbf{reward}（回报——"看护之外"）、\textbf{regard as}（视为：\textit{regard sth as important}）。"regardless"字面"无顾"——毫无顾虑。
\end{origin}

\vp{regardless of}

\sense{短语}{不管，不顾}
\ex{Everyone, regardless of age or background, is welcome to join the club.}
\ex{She chose to pursue her dream, regardless of the difficulties.}

\begin{gram}{regardless of / irrespective of / in spite of 的"不管"细分}
\textbf{regardless of}（不管——通用：\textit{regardless of the outcome}）；\textbf{irrespective of}（不考虑——更正式：\textit{irrespective of race} 不分种族）；\textbf{in spite of}（尽管——含"对抗"意味：\textit{in spite of the rain} 冒雨）。三者在"不顾条件"义上可互换，但 in spite of 强调"与困难对抗"的态度。
\end{gram}

\vw{subsequent}{/ˈsʌbsɪkwənt/}{adj.}

\sense{adj.}{随后的，后来的}
\col{subsequent to sth}
\ex{Subsequent to his resignation, the company went through a difficult period.}
\col{in subsequent years/months}
\ex{In subsequent years, the two countries became close allies.}

\begin{gram}{subsequent / following / next / later / eventual 的"随后"细分}
\textbf{subsequent}（随后的——正式：\textit{subsequent events} 后续事件）；\textbf{following}（随后的——次正式：\textit{the following day} 次日）；\textbf{next}（下一个——通用）；\textbf{later}（后来——时间后置）；\textbf{eventual}（最终的——强调结局：\textit{the eventual outcome}）。
\end{gram}

\begin{origin}
15 世纪借自拉丁语 \textit{subsequentem}（紧接的），是 \textit{subsequi}（紧随）的现在分词：\textit{sub-}（在下面）+ \textit{sequi}（跟随）。同源词 \textbf{subsequence}（随后发生之事）、\textbf{consequence}（结果；\textit{con-} 一起 + sequi $\to$ 随之而来的事物）、\textbf{sequence}（序列；连续）、\textbf{consequential}（重要的；自大的）、\textbf{execute}（执行；\textit{ex-} 出 + sequi $\to$ 跟随到底$\to$执行）、\textbf{prosecute}（起诉；\textit{pro-} 前 + sequi $\to$ 追到前面$\to$起诉）、\textbf{pursue}（追逐；\textit{pur-}（pro- 的变体）前 + sequi）、\textbf{ensue}（接着发生；\textit{en-} 入 + sequi $\to$ 跟随而入）。
\end{origin}

\vw{captain}{/ˈkæptɪn/}{n.}

\sense{n.}{船长，机长}
\col{the captain of a ship/plane}
\ex{The captain of the aircraft welcomed the passengers on board.}

\sense{n.}{上校}
\col{army/marine captain}
\ex{The captain led his troops into battle with great courage.}

\sense{n.}{队长}
\col{team captain 队长}
\ex{She was appointed team captain for the upcoming Olympic Games.}

\begin{gram}{captain / major / colonel / general 的军衔阶梯}
\textbf{captain}（上尉/船长/队长——校级最低）；\textbf{major}（少校）；\textbf{colonel}（上校）；\textbf{general}（将军）。注意 army captain 通常译"上尉"而非"上校"——captain 在不同军种译法不同。
\end{gram}

\begin{origin}
14 世纪借自古法语 \textit{capitaine}，拉丁语 \textit{capitaneus}（首领），源于 \textit{caput}（头）。同源词 \textbf{captain}、\textbf{capital}（首都；资本；\textit{caput} 头$\to$ 首要之物）、\textbf{cattle}（牛——拉丁语 \textit{capitale} 财产，原指"家畜头数"）、\textbf{chapter}（章节；\textit{caput} 头$\to$一本书的"头"）、\textbf{achieve}（达到；古法语 \textit{a chief} $\to$ \textit{a-} 到 + \textit{chef} 头$\to$ 到达顶端）、\textbf{chief}（首领；主要的）。
\end{origin}

\vp{allow for}

\sense{短语}{使成为可能；考虑到，把……计算在内}
\col{allow for sth}
\ex{The schedule allows for an extra hour of revision before the exam.}
\ex{The plan allows for unforeseen delays and additional costs.}

\begin{gram}{allow for / allow of / allow to do 的"允许"细分}
\textbf{allow for}（考虑到——预留空间：\textit{allow for delays} 预留延误）；\textbf{allow of}（容许——正式：\textit{The situation allows of no delay} 不容耽搁）；\textbf{allow sb to do sth}（允许某人做某事：\textit{allow children to play}）。三短语中只有 allow for 表"留余地"——预判未知。
\end{gram}

\vw{kingdom}{/ˈkɪŋdəm/}{n.}

\sense{n.}{王国}
\col{the United Kingdom 联合王国（英国正式名）}
\ex{The United Kingdom consists of England, Scotland, Wales, and Northern Ireland.}
\col{the animal/plant kingdom 动物/植物界}
\ex{The Amazon rainforest is home to a vast variety of species from the plant and animal kingdoms.}

\sense{n.}{管辖范围}
\col{within one's kingdom}
\ex{Within the kitchen, she is queen of her own kingdom.}

\begin{gram}{kingdom / realm / domain 的"领域"细分}
\textbf{kingdom}（王国；常用于 the animal/plant kingdom）；\textbf{realm}（王国；更文学：\textit{the realm of literature}）；\textbf{domain}（领域；偏抽象/学术：\textit{the domain of computer science}）；\textbf{territory}（领土；偏地理）。kingdom 偏生物/历史，realm 偏文学，domain 偏学术。
\end{gram}

\begin{origin}
古英语 \textit{cyingdōm}（王国），\textit{cying-}（王）+ \textit{-dōm}（状态；领域）。同源词 \textbf{king}（国王）、\textbf{kingly}（国王的）、\textbf{kingdom}、\textbf{dukedom}（公爵领地）、\textbf{earldom}（伯爵领地）、\textbf{realm}（王国；\textit{realme} 的法语形式）。-dom 后缀是古英语抽象名词后缀——\textit{freedom, wisdom, boredom} 皆此构造。
\end{origin}

\vw{astonish}{/əˈstɒnɪʃ/}{vt.}

\sense{vt.}{使十分惊讶，使吃惊}
\col{astonish sb}
\ex{The magician's tricks astonished the entire audience.}
\col{be astonished at/by sth}
\ex{I was astonished at/by the speed of his recovery.}

\begin{lex}{astonish / surprise / amaze / astound / startle 的"惊"强度谱}
\textbf{startle}（惊——轻且突然：\textit{startled by a noise}）；\textbf{surprise}（惊讶——通用：\textit{surprised at the news}）；\textbf{astonish}（震惊——较强：\textit{astonished at the courage}）；\textbf{amaze}（惊叹——更强：\textit{amazed at the view}）；\textbf{astound}（震惊至极——最强：\textit{astounded by the discovery}）。astonish 比 surprise 更强，amaze/astound 更戏剧化。
\end{lex}

\begin{origin}
15 世纪借自古法语 \textit{estoner}（使震惊），拉丁语 \textit{extonare}（打雷）：\textit{ex-}（出）+ \textit{tonare}（打雷）。同源词 \textbf{astonish}、\textbf{astonishment}（惊讶）、\textbf{astounding}（令人震惊的）、\textbf{astound}（使震惊；astound 是 astonish 的简写）、\textbf{ thunder }（雷；\textit{tonare} 英语化）、\textbf{thunderstruck}（震惊的；字面"被雷击"）。astonish 字面"被雷击中"——震惊如同雷击。
\end{origin}

\vw{royal}{/ˈrɔɪəl/}{adj.}

\sense{adj.}{盛大的，庄严的}
\col{a royal welcome}
\ex{The visiting head of state was given a royal welcome at the airport.}
\col{royal treatment}
\ex{The hotel guests received royal treatment during their stay.}

\sense{adj.}{皇家的，王室的}
\col{the royal family/household}
\ex{The royal family appears on the balcony of the palace on national holidays.}
\col{royal blood 皇室血统}
\ex{Princes are supposed to have royal blood, but the queen's youngest son is said to be a commoner.}

\begin{gram}{royal / regal / sovereign / imperial / majestic 的"王"细分}
\textbf{royal}（王室的；通用：\textit{royal family}）；\textbf{regal}（帝王的；偏文学/庄严：\textit{regal bearing} 帝王仪态）；\textbf{sovereign}（至高无上的；主权：\textit{sovereign state} 主权国家）；\textbf{imperial}（帝国的；\textit{the Imperial Palace} 皇宫）；\textbf{majestic}（雄伟庄严的；多形容自然/建筑：\textit{majestic mountains}）。royal 偏英国王室，imperial 偏历史帝国。
\end{gram}

\begin{origin}
14 世纪借自古法语 \textit{roial}，拉丁语 \textit{regalis}（王室的；属于王的），源于 \textit{rex}（王）。同源词 \textbf{royal}、\textbf{royalist}（保皇党人）、\textbf{regal}（帝王的）、\textbf{regalia}（王权的象征物；王冠、权杖）、\textbf{regent}（摄政者；\textit{regens} 统治者）、\textbf{regime}（政权；政体；\textit{regimen} 统治$\to$政治制度）、\textbf{region}（地区；\textit{regere} 统治$\to$被统治之地）、\textbf{regular}（规则的；\textit{regula} 尺子$\to$以尺子规则）、\textbf{rule}（规则；统治；\textit{regula} 的英语化）、\textbf{reign}（统治时期；\textit{regnum} 统治）。
\end{origin}

\vw{literary}{/ˈlɪtərəri/}{adj.}

\sense{adj.}{文学的，文学上的}
\col{literary criticism/works}
\ex{She is known for her contributions to modern literary criticism.}
\col{a literary prize}
\ex{The novel won the country's most prestigious literary prize.}

\sense{adj.}{爱好文学的}
\col{a literary circle}
\ex{He moved in literary circles in London during the 1920s.}

\begin{origin}
17 世纪借自拉丁语 \textit{literarius}（关于字母的；关于学问的），源于 \textit{litera}（字母）。同源词 \textbf{literary}、\textbf{literature}（文学；文献）、\textbf{literal}（字面的；\textit{litera} 字$\to$ 按字面）、\textbf{literate}（识字的；有文化的）、\textbf{illiterate}（不识字的；\textit{il-} 不 + \textit{litera}）、\textbf{alliteration}（头韵；\textit{al-} 朝 + \textit{litera} $\to$ 头字母相同）。"文学"字面"字母之学"——从识字到学问的递进。
\end{origin}

\vw{extensively}{/ɪkˈstensɪvli/}{adv.}

\sense{adv.}{广阔地；广泛地}
\col{extensively studied/researched}
\ex{The phenomenon has been extensively studied over the past two decades.}
\col{read/travel extensively}
\ex{She has travelled extensively in South America and Asia.}

\begin{gram}{extensively / widely / broadly / largely 的"广"细分}
\textbf{extensively}（广泛地——偏"规模大"：\textit{extensively researched} 大规模研究的）；\textbf{widely}（广泛地——偏"覆盖广"：\textit{widely read} 读者众多）；\textbf{broadly}（广义地——偏"范围广"：\textit{broadly speaking}）；\textbf{largely}（主要地——偏"占大部分"：\textit{largely due to}）。
\end{gram}

\begin{origin}
17 世纪由 \textit{extensive}（广泛的）+ \textit{-ly} 派生。extensive 借自拉丁语 \textit{extensivus}（延展的），源于 \textit{extendere}（向外伸展）。同源词 \textbf{extensive}（广泛的；广阔的）、\textbf{extend}（延伸）、\textbf{extent}（程度；范围）、\textbf{extension}（延伸；延期）、\textbf{tension}（张力；\textit{tendere} 拉）。"广泛"字面"伸展得远"——空间/范围上的延展。
\end{origin}

\vw{facilitate}{/fəˈsɪlɪteɪt/}{vt.}

\sense{vt.}{使便利，促进}
\col{facilitate sth}
\ex{The new software will facilitate communication between team members.}
\col{facilitate the process/progress}
\ex{Good lighting facilitates the learning process in the classroom.}

\begin{gram}{facilitate 严格不接 for/from/to sb}
\textbf{facilitate} 直接接宾语：\ding{55} \textit{facilitate to do}；\checkmark\ \textit{facilitate doing}。类似规则动词：\textit{hesitate, anticipate, appreciate, tolerate}。例：\textit{Modern tools facilitate \textbf{communication}}（促进交流）。
\end{gram}

\begin{origin}
17 世纪借自拉丁语 \textit{facilitatem}（容易、便利），是 \textit{facilis}（容易的）的抽象名词。同源词 \textbf{facility}（设施；才能；便利）、\textbf{facile}（容易的；流畅的；\textit{a facile answer} 流于表面的回答）、\textbf{facilitate}（使便利）、\textbf{difficult}（困难的；\textit{dis-} 不 + facilis）、\textbf{faculty}（能力；学院）。
\end{origin}

\vp{sort out}

\sense{短语}{理顺，整理}
\col{sort out sth}
\ex{We need to sort out the mess before the guests arrive.}
\col{sort out the problem/issue}
\ex{I'm trying to sort out the problem with my bank account.}

\sense{短语}{把……安排好}
\col{sort sb out}
\ex{She sorted out the travel arrangements for the whole team.}
\ex{I'll sort you out — you'll have your money by tomorrow.}

\begin{gram}{sort out / sort through / sort of 的"分类"细分}
\textbf{sort out}（整理；解决）；\textbf{sort through}（翻查以分类：\textit{sort through a pile of papers}）；\textbf{sort of}（有几分——口语："\textit{I sort of like him} 我有点喜欢他"）。三短语中 sort 含义统一为"分类"，语境不同则义不同。
\end{gram}

\vw{available}{/əˈveɪləbl/}{adj.}

\sense{adj.}{可获得的，可找到的}
\col{available to/for sb}
\ex{Tickets are available to the general public from next Monday.}
\col{make sth available}
\ex{The university is making more resources available to first-year students.}

\sense{adj.}{有空的}
\col{be available (to do sth)}
\ex{I'm available to help you this weekend if you need me.}

\begin{gram}{available 作表语时的位置与定语位置的微妙差异}
\textbf{available} 作定语时常后置：\textit{the only ticket available}（唯一可得的票——非 \textit{available ticket}，后者虽可但较生硬）。这是英语"形容词后置"的特殊规则——available 严格说是表语形容词，作定语时倾向后置。同类的"前形容词"：\textit{the best hotel available} / \textit{the seats available}。
\end{gram}

\begin{origin}
15 世纪借自古法语 \textit{available}，源于 \textit{a-}（ad- 的同化形式"+vail"有益）+ \textit{vail}（价值、有益），拉丁语 \textit{valere}（强、有价值）。同源词 \textbf{avail}（有用；惠及）、\textbf{available}（可获得的）、\textbf{prevail}（盛行；\textit{prae-} 极 + \textit{valere} $\to$ 极有力量）、\textbf{valid}（有效的；\textit{valere} $\to$ 有力量的$\to$有效的）、\textbf{valour}（勇武）、\textbf{value}（价值）、\textbf{evaluate}（评估）、\textbf{equivalent}（等价的）。
\end{origin}

\vw{exceptional}{/ɪkˈsepʃənl/}{adj.}

\sense{adj.}{杰出的，优秀的}
\col{exceptional talent/skill}
\ex{The young pianist showed exceptional talent from an early age.}
\col{an exceptional student/performance}
\ex{It was an exceptional performance by the entire cast.}

\sense{adj.}{异常的，罕见的}
\col{exceptional circumstances/conditions}
\ex{The match was cancelled due to exceptional weather conditions.}
\col{an exceptional case}
\ex{This is an exceptional case, and the usual rules do not apply.}

\begin{gram}{exceptional / exceptionable / exceptional / exception 的派生分叉}
\textbf{exception}（例外）、\textbf{exceptional}（杰出的；异常的——双义）、\textbf{exceptionable}（可反对的；引起异议的——多用于法律）、\textbf{unexceptionable}（无可非议的）、\textbf{unexceptional}（平凡的）。注意 exceptional 是"例外"的两面——可褒可贬，例外既可杰出也可异常。
\end{gram}

\begin{origin}
17 世纪由 \textit{exception}（例外）+ \textit{-al} 派生。exception 14 世纪借自古法语 \textit{exception}，拉丁语 \textit{exceptionem}（例外），是 \textit{excipere}（取出、排除）的派生形式：\textit{ex-}（出）+ \textit{capere}（抓、取）。同源词 \textbf{except}（除……外）、\textbf{exception}（例外）、\textbf{exceptional}（杰出的；异常的）、\textbf{accept}（接受；\textit{ac-} 朝 + capere $\to$ 朝向取$\to$接受）、\textbf{concept}（概念；\textit{con-} 一起 + capere $\to$ 一起抓住$\to$ 概念）、\textbf{intercept}（拦截）、\textbf{recipient}（接受者；\textit{re-} 回 + capere）。
\end{origin}

\vw{phrase}{/freɪz/}{n.}

\sense{n.}{成语，习语}
\col{a set phrase 固定搭配}
\ex{``By and large'' is a set phrase meaning ``generally speaking''.}

\sense{n.}{短语，词组}
\col{a noun/verb/adjective phrase}
\ex{``The man in the corner'' is a noun phrase.}
\col{turn a phrase}
\ex{She has a wonderful way of turning a phrase that makes her writing sparkle.}

\begin{lex}{phrase / idiom / expression / clause / sentence 的"短语"细分}
\textbf{phrase}（短语；不构成完整句：\textit{a noun phrase}）；\textbf{idiom}（习语；固定表达：\textit{under the weather}）；\textbf{expression}（表达；泛指：\textit{a useful expression}）；\textbf{clause}（从句；有主谓的短语：\textit{a subordinate clause}）；\textbf{sentence}（句子）。
\end{lex}

\begin{origin}
16 世纪借自拉丁语 \textit{phrasis}（说法；表达），希腊语 \textit{phrasis}（言语、表达），源于 \textit{phrazein}（说话、表达）。同源词 \textbf{phrase}（短语）、\textbf{paraphrase}（意译；\textit{para-} 旁 + \textit{phrase} $\to$ 在旁边说$\to$ 解释）、\textbf{periphrasis}（迂回说法；\textit{peri-} 周围 + \textit{phrase} $\to$ 绕着说$\to$ 迂回）。"短语"字面"说出来的一组词"。
\end{origin}

\vw{format}{/ˈfɔːmæt/}{n. \& vt.}

\sense{n.}{总体安排，设计}
\col{the format of sth}
\ex{The format of the conference has changed this year, with more workshops and fewer lectures.}

\sense{n.}{版式}
\col{a magazine/book format}
\ex{The dictionary is now available in both print and digital formats.}

\begin{gram}{format / formation / formula / formality 的"形"细分}
\textbf{format}（版式；安排——\textit{file format} 文件格式）；\textbf{formation}（形成；编队——\textit{flying formation} 飞行编队）；\textbf{formula}（公式；配方——\textit{chemical formula} 化学式）；\textbf{formality}（形式；手续——\textit{a mere formality}）。四者皆源自 \textit{forma}（形状），但分工明确。
\end{gram}

\begin{origin}
19 世纪借自法语 \textit{format}，拉丁语 \textit{formatus}（已形成的），是 \textit{formare}（塑造、形成）的过去分词，源于 \textit{forma}（形状、形式）。同源词 \textbf{form}（形式；形状）、\textbf{formal}（正式的；形式的）、\textbf{formation}（形成；编队）、\textbf{formula}（公式）、\textbf{reform}（改革；\textit{re-} 再 + forma $\to$ 再塑$\to$ 改造）、\textbf{transform}（转变；\textit{trans-} 跨越 + forma）、\textbf{conform}（遵守；\textit{con-} 一起 + forma $\to$ 同形$\to$遵守）、\textbf{inform}（通知；\textit{in-} 入 + forma $\to$ 注入形式$\to$ 通知）。
\end{origin}

\vw{adopt}{/əˈdɒpt/}{vt.}

\sense{vt.}{采用，采取}
\col{adopt a method/approach/policy}
\ex{The school has adopted a new approach to teaching mathematics.}
\col{adopt a style/attitude}
\ex{She has adopted a more positive attitude towards her work.}

\sense{vt.}{表决采纳}
\col{adopt a resolution/law/constitution}
\ex{The committee adopted the new regulations without dissent.}

\sense{vt. \& vi.}{领养，收养}
\col{adopt a child}
\ex{They adopted a baby girl from China last year.}

\begin{gram}{adopt / adapt / adept 的形近义远}
\textbf{adopt}（采用；领养；\textit{ad-} 朝向 + \textit{optare} 选择）；\textbf{adapt}（适应；改编；\textit{ad-} 朝 + \textit{aptus} 适合）；\textbf{adept}（熟练的；\textit{ad-} 朝 + \textit{aptus} $\to$ 达到熟练的）。三词仅一字母之差但词义无关。
\end{gram}

\begin{origin}
16 世纪借自法语 \textit{adopter}，拉丁语 \textit{adoptare}（挑选、收养）：\textit{ad-}（朝向、为）+ \textit{optare}（选择、希望）。同源词 \textbf{adopt}、\textbf{adoption}（采纳；领养）、\textbf{adoptive}（收养的；\textit{adoptive parents} 养父母）、\textbf{option}（选择）、\textbf{optional}（可选的）、\textbf{co-opt}（吸收为成员）。
\end{origin}

\vw{comparison}{/kəmˈpærɪsn/}{n.}

\sense{n.}{对比；比较}
\col{in comparison with sth}
\ex{In comparison with last year, sales have increased by 15 per cent.}
\col{by comparison 相比之下}
\ex{By comparison, the second novel is much weaker than the first.}
\col{draw a comparison between A and B}
\ex{The report draws a comparison between the two education systems.}

\begin{gram}{comparison / contrast 的"比较"细分}
\textbf{comparison}（比较——找相似：\textit{a comparison of two methods}）；\textbf{contrast}（对比——找差异：\textit{a contrast between two styles}）。考试常考：\textit{by comparison} 相比之下找相似；\textit{by contrast} 相比之下找差异。
\end{gram}

\begin{origin}
15 世纪借自古法语 \textit{comparaison}，拉丁语 \textit{comparisonem}（比较），源于 \textit{comparare}（放在一起比较）：\textit{com-}（一起）+ \textit{par}（平等）。同源词 \textbf{compare}（比较）、\textbf{comparable}（可比较的）、\textbf{comparative}（比较的；相当的——\textit{comparative literature} 比较文学）、\textbf{par}（同等）、\textbf{parity}（平等）、\textbf{disparity}（不一致）、\textbf{impartial}（公正的；\textit{im-} 不 + pars $\to$ 不偏袒的）。
\end{origin}

\vw{pack}{/pæk/}{vt. \& vi. \& n.}

\sense{vt. \& vi.}{塞进，挤进}
\col{pack (sth) into sth}
\ex{She managed to pack all her clothes into a single suitcase.}
\ex{Thousands of fans packed into the stadium for the concert.}

\sense{vt. \& vi.}{装（箱），收拾；打包}
\col{pack (a bag/suitcase)}
\ex{Have you packed your bags yet? We're leaving in an hour.}
\col{pack for a trip}
\ex{I'm not very good at packing for long trips.}

\sense{vt. \& vi.}{包装}
\col{pack sth in/with sth}
\ex{The meat is packed in vacuum-sealed plastic to keep it fresh.}

\sense{n.}{包，包裹}
\col{a pack of sth}
\ex{She bought a pack of postcards at the gift shop.}

\sense{n.}{小盒，小包}
\col{a pack/packet of cigarettes}
\ex{He smoked a pack of cigarettes a day for twenty years.}

\begin{gram}{pack / packet / package / parcel / bundle 的"包"细分}
见 parcel 条。\textbf{pack}（包；通用；\textit{a pack of cards} 一副牌；\textit{a pack of cigarettes}）；\textbf{packet}（小包；\textit{a packet of crisps} 一袋薯片）；\textbf{package}（包裹；套装；\textit{a software package}）；\textbf{parcel}（邮寄包裹；\textit{a parcel post}）；\textbf{bundle}（捆扎）。
\end{gram}

\begin{origin}
13 世纪中古英语 \textit{pak}（包裹；捆），可能借自中世纪荷兰语 \textit{pak}（包裹）。同源词 \textbf{packet}（小包）、\textbf{packing}（包装）、\textbf{pact}（公约——不同源；拉丁语 \textit{pactum}），\textbf{pack} 作"群"（\textit{a pack of wolves}）、\textbf{package}（包裹）、\textbf{backpack}（背包）、\textbf{packed}（拥挤的；\textit{a packed train}）。
\end{origin}

\vw{contrary}{/ˈkɒntrəri/}{adj. \& n.}

\sense{adj.}{相反的，相对立的}
\col{contrary to sth}
\ex{Contrary to popular belief, the Great Wall is not visible from the moon with the naked eye.}
\col{contrary views/opinions}
\ex{The two politicians hold contrary views on almost every issue.}

\sense{adj.}{截然不同的}
\col{contrary to expectations}
\ex{Contrary to expectations, the meeting ended earlier than planned.}

\sense{n.}{相反（的事物）}
\col{the contrary}
\ex{Some people like the heat, but I prefer the contrary.}

\begin{lex}{contrary / opposite / contradictory / converse 的"反"细分}
\textbf{contrary}（相反的——多修饰观点/信念：\textit{contrary opinions}）；\textbf{opposite}（对面的；相反的——多修饰方向/位置：\textit{the opposite side}）；\textbf{contradictory}（矛盾的——多指逻辑上不可同时成立：\textit{contradictory statements}）；\textbf{converse}（逆命题；反面——多用于逻辑/哲学：\textit{the converse is also true}）。
\end{lex}

\begin{origin}
14 世纪借自古法语 \textit{contrarie}，拉丁语 \textit{contrarius}（相反的、相对立的），源于 \textit{contra}（反对、相对）。同源词 \textbf{contrary}、\textbf{contrary to}、\textbf{on the contrary}（正相反——\textit{I'm not bored; on the contrary, I'm having a great time}）、\textbf{contrast}（对比）、\textbf{encounter}（遭遇）、\textbf{contra}（反对）。
\end{origin}

\vp{contrary to}

\sense{短语}{与……相反}
\ex{Contrary to what he said, the meeting was very productive.}
\ex{Contrary to the weather forecast, it stayed sunny all day.}

\begin{gram}{contrary to / in contrast to / on the contrary / on the other hand 的"对照"短语}
\textbf{contrary to}（与……相反——用于事实：\textit{contrary to the report}）；\textbf{in contrast to}（与……对比——强调差异：\textit{in contrast to last year}）；\textbf{on the contrary}（正相反——用于反驳：\textit{``You're not tired?'' ``On the contrary, I'm exhausted.''}）；\textbf{on the other hand}（另一方面——用于平衡：\textit{on the other hand, we should consider...}）。
\end{gram}

\vw{hell}{/hel/}{n.}

\sense{n.}{地狱}
\col{in hell}
\ex{The medieval paintings of hell are terrifying images of eternal suffering.}

\sense{n.}{苦难的经历，悲惨的境况}
\col{a hell of ... (口语)}
\ex{The journey was a hell of a trip — five days of non-stop travel.}
\col{give sb hell 让某人受罪}
\ex{The boss gave him hell for being late again.}

\begin{gram}{the hell 作强调副词：what/why/where the hell}
\textbf{what the hell}（到底什么——\textit{What the hell are you doing here?}）、\textbf{why the hell}（到底为什么）、\textbf{where the hell}（到底哪里）。口语中加 the hell 表"强调/愤怒"——比单纯 what 强烈得多。\textbf{like hell}（拼命地；\textit{We worked like hell to finish on time}）。
\end{gram}

\begin{origin}
古英语 \textit{hel}（阴间；死后之境），原始日耳曼语 *\textit{haljō}（隐藏之所），与 \textit{helan}（隐藏）同源。同源词 \textbf{hell}、\textbf{helmet}（头盔——\textit{helan} 隐藏$\to$ 隐藏头部之物）、德语 \textit{Hölle}（地狱）、\textbf{conceal}（隐藏；\textit{con-} 一起 + \textit{celare} 隐藏；与 hell 形近但实不同源）。"地狱"原义"被隐藏之所"——死者的"看不见"之地。
\end{origin}

\vw{rub}{/rʌb/}{vt. \& vi.}

\sense{vt. \& vi.}{擦，磨，搓}
\col{rub sth (with/on sth)}
\ex{She rubbed her hands together to warm them up.}
\col{rub sth off/out 擦掉}
\ex{The ink mark on her shirt wouldn't rub out.}

\sense{vi.}{(使) 相互摩擦}
\col{rub against sth}
\ex{The cat rubbed against my leg as I came into the room.}

\begin{gram}{rub it in / rub salt in the wound：伤口撒盐的"擦"隐喻}
\textbf{rub it in}（反复提及令人不快的事；\textit{Stop rubbing it in — I know I made a mistake}）、\textbf{rub salt in the wound}（伤口上撒盐；雪上加霜）、\textbf{rub sb up the wrong way}（触犯某人；惹恼某人）、\textbf{rub shoulders with}（与名人接触）。
\end{gram}

\begin{origin}
14 世纪中古英语 \textit{rubben}，可能与古英语 \textit{ribbing}（摩擦）相关。同源词 \textbf{rub}、\textbf{rubber}（橡胶；橡皮——\textit{rub} + \textit{-er} $\to$ 用来擦的东西）、\textbf{rubbish}（垃圾；可能与 \textit{rub} 同源——"被擦下来的碎屑"）。rubber 的"橡皮"义源自其原本功能——擦掉铅笔字。
\end{origin}

\vw{uniform}{/ˈjuːnɪfɔːm/}{n. \& adj.}

\sense{n.}{制服；队服}
\col{in uniform}
\ex{The soldiers were all in uniform at the ceremony.}
\col{school uniform 校服}
\ex{All students are required to wear school uniform.}

\sense{adj.}{一致的，统一的}
\col{uniform colour/size}
\ex{The houses in the new estate are all of a uniform size and colour.}
\col{at a uniform speed/rate}
\ex{The machine operates at a uniform speed.}

\begin{lex}{uniform / consistent / constant / identical 的"一致"细分}
\textbf{uniform}（一致的——多修饰物理特性：\textit{uniform temperature}）；\textbf{consistent}（始终如一的——多修饰行为/态度：\textit{consistent policy}）；\textbf{constant}（恒定的——多修饰时间：\textit{constant speed}）；\textbf{identical}（完全相同的——多修饰可数事物：\textit{identical twins}）。uniform 偏物理，consistent 偏行为。
\end{lex}

\begin{origin}
16 世纪借自拉丁语 \textit{uniformis}（只有一种形状的），\textit{uni-}（一、单一）+ \textit{forma}（形状）。同源词 \textbf{uniform}、\textbf{uniformity}（一致性；一律）、\textbf{unify}（统一；使一致）、\textbf{union}（联合；工会）、\textbf{unique}（独特的；\textit{uni-} 一$\to$ 独一无二的）、\textbf{unit}（单位；个体）、\textbf{universe}（宇宙）、\textbf{universal}（普遍的）。
\end{origin}

\vw{violently}{/ˈvaɪələntli/}{adv.}

\sense{adv.}{猛烈地，厉害地}
\ex{The car crashed violently into the wall.}
\ex{The door shook violently in the strong wind.}

\sense{adv.}{强烈地，激烈地}
\ex{She reacted violently to the suggestion that she should retire.}

\sense{adv.}{凶狠地}
\ex{The suspect was known to behave violently towards his family.}

\begin{gram}{violently / fiercely / intensely / vigorously 的"猛烈"细分}
\textbf{violently}（猛烈地——含暴力倾向：\textit{react violently}）；\textbf{fiercely}（凶狠地；激烈地——\textit{fiercely competitive} 竞争激烈）；\textbf{intensely}（强烈地——偏情感/感官：\textit{intensely painful}）；\textbf{vigorously}（精力充沛地——偏活动：\textit{vigorously denied} 坚决否认）。violently 偏"含暴力可能"，fiercely 偏"竞争性激烈"。
\end{gram}

\begin{origin}
15 世纪由 \textit{violent}（暴力的；猛烈的）+ \textit{-ly} 派生。violent 14 世纪借自古法语 \textit{violent}，拉丁语 \textit{violentus}（猛烈的、暴力的），源于 \textit{vis}（力量）。同源词 \textbf{violent}（暴力的）、\textbf{violence}（暴力）、\textbf{violate}（违反；侵犯；\textit{vis} 力量 $\to$ 强用力量$\to$侵犯）、\textbf{violin}（小提琴；与 violent 形近但实不同源——意大利语 \textit{violino} 小型中提琴）、\textbf{virile}（有男子气概的；\textit{virilis} 男人的）。
\end{origin}

\vw{dive}{/daɪv/}{vi. \& n.}

\sense{vi.}{俯冲}
\col{dive into/towards sth}
\ex{The hawk dived towards its prey.}
\col{dive-bomb 俯冲轰炸}
\ex{The aircraft dived-bombed the enemy positions.}

\sense{vi.}{跳水}
\col{dive into the water}
\ex{She dived into the pool from the highest board.}

\sense{vi.}{潜水}
\col{dive for pearls/shells}
\ex{The fishermen dive for pearls off the coast of Japan.}

\sense{n.}{冲，扑，俯冲}
\col{take a dive}
\ex{The goalkeeper took a dive in the penalty area — or did he?}

\sense{n.}{潜水；跳水}
\col{a deep dive 深入探究（习语）}
\ex{The report takes a deep dive into the causes of the financial crisis.}

\begin{gram}{dive / plunge / plummet / dip 的"跌"细分}
\textbf{dive}（俯冲——快速下降：\textit{the bird dived}）；\textbf{plunge}（跳入——猛然下坠：\textit{the price plunged}）；\textbf{plummet}（暴跌——垂直下落：\textit{stocks plummeted}）；\textbf{dip}（轻降——温和下降：\textit{temperatures dipped}）。dive 偏快速且有意，plunge 偏突然，plummet 偏指数级，dip 偏轻微。
\end{gram}

\begin{origin}
古英语 \textit{dȳfan}（使下沉；浸入），与 \textit{dūfe}（潜水鸟）同源，原始日耳曼语 *\textit{dūbjan}。同源词 \textbf{dive}、\textbf{diver}（潜水员）、\textbf{diving board}（跳板）。与 \textit{dove}（鸽子）形近但实不同源——dove 来自古英语 \textit{dūfe} 潜水鸟。
\end{origin}

\vw{freeze}{/friːz/}{vi. \& vt.}

\sense{vi.}{突然停止，呆住}
\col{freeze in one's tracks}
\ex{The deer froze in its tracks when it saw the hunter.}
\col{freeze with fear/terror}
\ex{She froze with terror when she heard the strange noise.}

\sense{vi.}{结冰}
\col{freeze over/up}
\ex{The lake freezes over every winter.}

\sense{vi.}{冻僵}
\ex{Two climbers were rescued after they froze to death on the mountain.}

\begin{gram}{freeze / refrigerate / chill / cool 的"冷"细分}
\textbf{freeze}（冻——水变冰：\textit{freeze the meat}）；\textbf{refrigerate}（冷藏——0 度以上保存：\textit{refrigerate the milk}）；\textbf{chill}（冷却——降温但未必结冰：\textit{chill the wine}）；\textbf{cool}（冷却——泛指降温：\textit{cool the soup}）。freeze 是 0°C 以下结冰，chill 是冷藏温度，cool 是温和降温。
\end{gram}

\begin{origin}
古英语 \textit{frēosan}（冻结、冻死），原始日耳曼语 *\textit{frus-}（结冰）。同源词 \textbf{freeze}、\textbf{freezer}（冷冻柜）、\textbf{freezing}（冰冻的）、\textbf{frost}（霜；\textit{freeze} 的派生）、\textbf{frozen}（结冰的）、德语 \textit{frieren}（冻）。"freeze" 与 "frost" 形近同根——frost 即"已结成的冰"。
\end{origin}

\vw{deck}{/dek/}{n.}

\sense{n.}{甲板，舱面}
\col{on deck}
\ex{The passengers were sunbathing on deck.}
\col{below deck(s) 在甲板下}
\ex{Crew members are not allowed to socialise with passengers below decks.}

\begin{gram}{deck / floor / ground / level 的"层"细分}
\textbf{deck}（甲板——船/飞机/建筑物的某一层；\textit{the upper deck of a bus} 双层巴士的上层）；\textbf{floor}（地板——建筑物的层；\textit{the second floor}）；\textbf{ground}（地面）；\textbf{level}（层次；\textit{on a higher level}）。deck 多用于船/飞机/双层巴士。
\end{gram}

\begin{origin}
15 世纪中古英语 \textit{dekke}（甲板），可能与中世纪荷兰语 \textit{dek}（覆盖物；甲板）同源——\textit{deck} 原义"覆盖在船上的板"。同源词 \textbf{deck}、\textbf{deckchair}（折叠躺椅——在甲板上用的椅子）。
\end{origin}

\vw{meanwhile}{/ˈmiːnwaɪl/}{adv.}

\sense{adv.}{同时}
\ex{Mother was preparing dinner. Meanwhile, the children were playing in the garden.}

\sense{adv.}{其间}
\ex{The doctor will be here in ten minutes. In the meanwhile, please try to keep the patient calm.}

\sense{adv.}{对比之下}
\ex{The company's profits have fallen sharply. Meanwhile, its competitors continue to thrive.}

\begin{gram}{meanwhile / meantime / in the meantime 的"其间"细分}
\textbf{meanwhile}（同时；其间——偏正式，连接两个事件）；\textbf{meantime}（其间——多用于 \textit{in the meantime}）；\textbf{in the meantime}（与此同时；在此期间）。三者在"期间"义上可互换，但 meanwhile 更适合做连接副词（\textit{Meanwhile, the war continued}）。
\end{gram}

\begin{origin}
14 世纪由 \textit{mean}（中间）+ \textit{while}（时间）合成。同源词 \textbf{mean}（中间的；\textit{the mean value} 平均值）、\textbf{while}（一段时间）。"与此同时"字面"在中间的时间里"。
\end{origin}

\vw{drag}{/dræɡ/}{vt. \& vi.}

\sense{vt.}{拖，拉}
\col{drag sth (across/along/over sth)}
\ex{He dragged the heavy suitcase across the floor.}
\col{drag sb away/back/out}
{The guards dragged the prisoner away.}

\sense{vi.}{缓慢而费力地移动}
\col{drag on 拖拖拉拉}
\ex{The meeting dragged on for over three hours with no decision reached.}

\begin{gram}{drag on / drag out / drag in 的"拖"三态}
\textbf{drag on}（拖长：\textit{The war dragged on for years}）；\textbf{drag out}（使拖长：\textit{drag out the negotiations} 拖延谈判）；\textbf{drag in}（硬扯进来：\textit{drag in irrelevant details} 扯进无关细节）。三短语中"on/out/in"皆为方向性小品词。
\end{gram}

\begin{origin}
15 世纪中古英语 \textit{draggen}，可能与古英语 \textit{dragan}（拖、拉）同源，原始日耳欧语 *\textit{drag-}。同源词 \textbf{dragon}（龙——\textit{dragan} 拉 $\to$ 拉长的怪物$\to$龙）、\textbf{dragnet}（拖网）、\textbf{dray}（板车；载重车）、德语 \textit{tragen}（承担）。"drag" 与 "dragon" 同根——古英语中 dragon（龙）是"被拉长之物"。
\end{origin}

\vw{chaos}{/ˈkeɪɒs/}{n.}

\sense{n.}{混乱，杂乱}
\col{in chaos}
\ex{The country has been in chaos since the government collapsed.}
\col{complete/total chaos}
\ex{After the party, the house was in complete chaos.}
\col{descend into chaos}
\ex{The demonstration quickly descended into chaos.}

\begin{origin}
15 世纪借自拉丁语 \textit{chaos}，希腊语 \textit{chaos}（混沌、无序的虚空），源于 \textit{chasma}（裂口、张开的口）、\textit{chaskein}（打哈欠、张开）。同源词 \textbf{chaotic}（混乱的）、\textbf{chaotically}（混乱地）、\textbf{chasm}（深渊；裂口；差距——\textit{a chasm between the rich and the poor} 贫富鸿沟）。希腊神话中 chaos 是宇宙诞生前的"虚空"——一切未分。
\end{origin}

\vw{hatch}{/hætʃ/}{vt. \& vi.}

\sense{vt.}{密谋，策划}
\col{hatch a plan/plot/scheme}
\ex{They hatched a plan to escape from the prison.}
\col{hatch out 密谋出}
\ex{The rebels hatched out a daring plot to overthrow the government.}

\sense{vt.}{使孵出}
\col{hatch (out) chicks/eggs}
\ex{The hen hatched out a dozen chicks.}

\sense{vi.}{孵化；孵出}
\col{hatch (out)}
\ex{The eggs take about three weeks to hatch.}

\begin{gram}{hatch 作名词"舱口"的另一义}
\textbf{hatch}（n. 舱口；厨房传菜窗口）：\textit{escape hatch} 应急出口、\textit{a serving hatch} 传菜口。同形异义：hatch（孵化）源自古英语 \textit{hæccan}，hatch（舱口）源自古英语 \textit{hæcc}（门闩）。
\end{gram}

\begin{origin}
13 世纪中古英语 \textit{hacchen}（孵出），可能拟声于"啄壳声"——小鸡破壳的"咔嗒"声。同源词 \textbf{hatchery}（孵卵处）、\textbf{hatchback}（掀背式汽车——后门"掀开"的方式与 hatch 同形不同源）。"hatch a plot"是"秘密策划"——计划像鸡蛋一样"暗中孵出"。
\end{origin}

\vw{plot}{/plɒt/}{n. \& vi. \& vt.}

\sense{n.}{阴谋}
\col{a plot to do sth}
\ex{The police uncovered a plot to assassinate the president.}
\col{foil/thwart a plot}
\ex{The spy managed to foil the enemy's plot.}

\sense{n.}{故事情节}
\col{the plot of a novel/play}
\ex{The plot of the novel is set in 19th-century Paris.}
\col{plot twist 情节转折}
\ex{The film has so many plot twists that it's hard to follow.}

\sense{vi. \& vt.}{密谋}
\col{plot against sb}
\ex{They plotted against the king and were eventually caught.}
\col{plot to do sth}
\ex{The rebels were plotting to overthrow the government.}

\begin{lex}{plot / scheme / conspire / intrigue 的"密谋"细分}
\textbf{plot}（密谋——多指小规模：\textit{plot a crime}）；\textbf{scheme}（策划——偏中性，可褒可贬：\textit{scheme a way out} 策划出路；\textit{tax scheme}）；\textbf{conspire}（共谋——多指多人联合：\textit{conspire to commit fraud}）；\textbf{intrigue}（密谋——偏戏剧性：\textit{political intrigue} 政治阴谋）。plot 偏密谋具体行为，scheme 偏计划/方案。
\end{lex}

\begin{origin}
古英语 \textit{plot}（一小块地），与 \textit{plot}（阴谋；情节）实为不同源——情节义 16 世纪起出现，可能是 \textit{complot}（共谋）简化形式。同源词 \textbf{plot}（一小块地；密谋；情节）、\textbf{plotter}（密谋者）、\textbf{grass plot}（草地）。注意 plot 的三义共形但来源不同。
\end{origin}

\vw{infamy}{/ˈɪnfəmi/}{n.}

\sense{n.}{臭名昭著，声名狼藉}
\col{eternal/everlasting infamy}
\ex{The dictator's name is remembered in eternal infamy.}
\col{a day of infamy 耻辱日}
\ex{President Roosevelt called the Japanese attack on Pearl Harbor ``a date which will live in infamy''.}

\sense{n.}{恶行，罪恶}
\col{the infamies of war}
\ex{The novel describes the infamies of war in graphic detail.}

\begin{gram}{infamy / fame / defame 的词义光谱}
\textbf{infamy}（臭名——完全负面的名声）；\textbf{fame}（名声——中性，可好可坏：\textit{fame and fortune} 名与利）；\textbf{defame}（诽谤——动词：损害他人名誉）；\textbf{infamous}（声名狼藉的——形容词）。infamy 比 fame 强度更高，附加"道德谴责"的含义。
\end{gram}

\begin{origin}
15 世纪借自古法语 \textit{infamie}，拉丁语 \textit{infamia}（坏名声、耻辱），源于 \textit{in-}（不）+ \textit{fama}（名声、传闻）。同源词 \textbf{infamous}（声名狼藉的）、\textbf{defame}（诽谤）、\textbf{fame}（名声）、\textbf{famous}（著名的）、\textbf{famed}（出名的）、\textbf{rumour}（流言；与 \textit{fama} 同源——\textit{fama} 传闻$\to$ \textit{rumor}）。"infamy" 字面"无名声"——失去好名声后的"臭名"。
\end{origin}

\vw{sink}{/sɪŋk/}{vi. \& vt.}

\sense{vi.}{下沉，沉没}
\col{sink (like a stone)}
\ex{The ship sank within minutes of hitting the iceberg.}
\col{sink or swim 沉或浮（成败在此一举）}
\ex{When you start a new business, it's sink or swim for the first few months.}

\sense{vi.}{坐下}
\col{sink into a chair/the sofa}
\ex{She sank into the armchair with a sigh of relief.}

\sense{vi.}{降低}
\col{sink (to a low level)}
\ex{Temperatures will sink below zero tonight.}
\col{his heart sank 他的心沉了下去}
\ex{My heart sank when I saw the long queue.}

\begin{gram}{sink / drown / submerge / plunge 的"沉"细分}
\textbf{sink}（沉——自然下沉：\textit{the ship sank}）；\textbf{drown}（溺水——人/动物被水淹：\textit{he drowned in the river}）；\textbf{submerge}（浸没——主动或被动沉入水中：\textit{submerge the vegetables}）；\textbf{plunge}（猛投——快速投入：\textit{plunge into the water}）。sink 多指船只/物体，drown 多人/动物。
\end{gram}

\begin{origin}
古英语 \textit{sincan}（沉、下沉），原始日耳曼语 *\textit{sengwan}。同源词 \textbf{sink}、\textbf{sinking}（下沉；沉没）、德语 \textit{sinken}（沉）、\textbf{shrink}（收缩；\textit{sincan} 沉 $\to$ 萎缩）、古英语 \textit{sincan} 派生的 \textbf{sag}（松弛；下垂）。sink 与 swim 是英语"沉浮"的标准对偶。
\end{origin}

\vw{memorial}{/məˈmɔːriəl/}{n. \& adj.}

\sense{n.}{纪念碑（或像等）；纪念物}
\col{a war memorial 战争纪念碑}
\ex{A war memorial has been erected in the centre of the town.}
\col{a memorial to sb/sth}
\ex{The statue is a memorial to the victims of the disaster.}

\begin{origin}
15 世纪借自拉丁语 \textit{memoriale}（纪念物），是 \textit{memoria}（记忆）的派生形式，源于 \textit{memor}（记在心里的）。同源词 \textbf{memory}（记忆）、\textbf{memorable}（难忘的）、\textbf{memorise}（记住）、\textbf{remember}（记得；\textit{re-} 再 + \textit{memor} $\to$ 再记起）、\textbf{remembrance}（纪念；回忆）、\textbf{commemorate}（纪念；\textit{com-} 一起 + memorare $\to$ 一起记起$\to$ 共同纪念）。
\end{origin}

\vw{pray}{/preɪ/}{vi. \& vt.}

\sense{vi. \& vt.}{企盼，祈求}
\col{pray for sth}
\ex{We prayed for peace in the troubled region.}
\col{pray (to sb) to do sth}
\ex{She prayed to God to give her strength.}

\sense{vi. \& vt.}{祈祷，祷告}
\col{pray to God/the saints}
\ex{The family knelt down and prayed together.}

\begin{gram}{pray / hope / wish / beg 的"愿"细分}
\textbf{pray}（祈祷——多含宗教意涵：\textit{pray for forgiveness}）；\textbf{hope}（希望——常态：\textit{I hope so}）；\textbf{wish}（祝愿；难以实现的希望：\textit{I wish I were taller}）；\textbf{beg}（乞求——低姿态：\textit{beg for mercy}）。pray 比 hope 更"虔诚"，hope 是日常用语。
\end{gram}

\begin{origin}
13 世纪借自古法语 \textit{preier}，拉丁语 \textit{precari}（恳求、祈祷），源于 \textit{prex}（请求、祈祷）。同源词 \textbf{prayer}（祈祷）、\textbf{pray}、\textbf{precarious}（不确定的；危险的；\textit{prex} 请求 $\to$ 通过请求得到的$\to$ 不牢靠的）、\textbf{deprecate}（反对；轻视；\textit{de-} 下 + \textit{precari} $\to$ 祈祷避邪$\to$ 反对）、\textbf{imprecate}（诅咒；\textit{im-} 入 + \textit{precari} $\to$ 召唤不祥）。
\end{origin}

\vw{episode}{/ˈepɪsəʊd/}{n.}

\sense{n.}{一段经历，片段}
\col{an episode of ...}
\ex{That episode during the war was one of the darkest in our family's history.}
\col{a brief episode}
\ex{The two leaders met for a brief episode of negotiations before tensions flared again.}

\sense{n.}{一集}
\col{the final/next episode}
\ex{Don't miss the final episode of the series, airing this Sunday.}
\col{an episode of a TV show/series}
\ex{The show has been praised for its compelling episodes.}

\begin{gram}{episode / incident / event / occurrence 的"事"细分}
\textbf{episode}（一段经历——常带戏剧性或可叙述：\textit{an episode in his life}）；\textbf{incident}（事件——多指不寻常/负面的：\textit{a diplomatic incident} 外交事件）；\textbf{event}（事件——通用：\textit{historical events}）；\textbf{occurrence}（发生；偶发事件——偏中性）。
\end{gram}

\begin{origin}
17 世纪借自希腊语 \textit{epeisodion}（插入部分；进入部分），源于 \textit{eisodos}（进入）：\textit{eis}（入）+ \textit{hodos}（路）。同源词 \textbf{episode}、\textbf{episodic}（偶发的；断续的）、\textbf{period}（时期；\textit{peri-} 周围 + \textit{hodos} $\to$ 绕行一周的路程）。"episode" 字面"进入原剧"——希腊戏剧中场间插入的小故事。
\end{origin}

\vw{textile}{/ˈtekstaɪl/}{n.}

\sense{n.}{纺织品；(textiles) 纺织业}
\col{the textile industry}
\ex{The textile industry is one of the largest employers in the region.}
\col{a textile factory/mill 纺织厂}
\ex{Textile factories once lined the river, taking advantage of the water power.}

\begin{origin}
17 世纪借自拉丁语 \textit{textilis}（编织的；织成的），是 \textit{texere}（编织）的形容词。同源词 \textbf{textile}、\textbf{text}（原文；正文——\textit{texere} 编织$\to$ 编织之物$\to$ 文本）、\textbf{texture}（质地；肌理）、\textbf{context}（上下文；\textit{con-} 一起 + \textit{texere} $\to$ 编织在一起$\to$语境）、\textbf{pretext}（借口；\textit{prae-} 前 + \textit{texere} $\to$ 在前面编造$\to$ 借口）、\textbf{subtle}（微妙的；\textit{sub-} 下 + \textit{texere} $\to$ 在下面编织$\to$ 难以察觉的）。
\end{origin}

\vw{reform}{/rɪˈfɔːm/}{vt. \& n.}

\sense{vt.}{改进，改良，改革}
\col{reform sth}
\ex{The government is planning to reform the tax system.}
\col{reform oneself}
\ex{After years of bad behaviour, he finally decided to reform himself.}

\sense{n.}{改革，改良}
\col{education/health/electoral reform}
\ex{The new law is part of a wider programme of social reform.}
\col{reforms (pl.)}
\ex{The new government has introduced sweeping reforms across the public sector.}

\begin{gram}{reform / transform / deform / conform 的"形"派生}
\textbf{re-}（再）= 改好；\textbf{trans-}（跨越）= 形态改变；\textbf{de-}（去除）= 失去形状；\textbf{con-}（一起）= 同形$\to$ 遵守。reform 偏"改善"，transform 偏"完全改变"（\textit{butterfly transformation} 蝴蝶变形），deform 偏"破坏形状"，conform 偏"遵守形式"。
\end{gram}

\begin{origin}
14 世纪借自古法语 \textit{reformer}，拉丁语 \textit{reformare}（再形成、改造）：\textit{re-}（再）+ \textit{formare}（形成）。同源词 \textbf{reform}、\textbf{reformation}（改革）、\textbf{reformer}（改革者）、\textbf{reformatory}（少管所；改革场所）、\textbf{transform}（转变）、\textbf{deform}（变形）、\textbf{conform}（遵守）、\textbf{uniform}（一致的）。
\end{origin}

\vw{output}{/ˈaʊtpʊt/}{n.}

\sense{n.}{产量，输出量}
\col{industrial/economic output}
\ex{Industrial output rose by 4 per cent last year.}
\col{output per worker 人均产出}
\ex{The factory's output per worker has increased significantly over the past decade.}

\begin{lex}{output / outcome / outlet / outlook 的"out-"派生分叉}
\textbf{output}（产量；输出——\textit{industrial output}）；\textbf{outcome}（结果；\textit{the outcome of the election}）；\textbf{outlet}（出口；发泄渠道：\textit{an outlet for one's emotions}）；\textbf{outlook}（前景；观点：\textit{a positive outlook}）。皆由 \textit{out-} 派生，但引申义各异。
\end{lex}

\begin{origin}
20 世纪初由 \textit{out-}（外、出）+ \textit{put}（放置）合成。同源词 \textbf{output}、\textbf{input}（投入；输入）、\textbf{throughput}（吞吐量）、\textbf{put}（放；表达）、\textbf{out}（外面）、\textbf{putout}（出局）。output 与 input 配对——产出与投入是现代经济学的核心对偶。
\end{origin}

\vw{shift}{/ʃɪft/}{n. \& vi. \& vt.}

\sense{n.}{改变，转移}
\col{a shift in sth}
\ex{There has been a major shift in public opinion on the issue.}
\col{a shift of emphasis/attention}
\ex{The new management brought about a shift of emphasis from profit to quality.}

\sense{n.}{轮班}
\col{work the day/night shift}
\ex{She works the night shift at the hospital.}
\col{an eight-hour shift}
\ex{Workers are demanding shorter shifts and better pay.}

\sense{vi. \& vt.}{(使) 移动，(使) 转移}
\col{shift (sth) from ... to ...}
\ex{She shifted her weight from one foot to the other.}
\col{shift gears 换挡}
\ex{The driver shifted gears and accelerated onto the motorway.}

\begin{gram}{shift / move / transfer / relocate 的"移"细分}
\textbf{shift}（转移——多指抽象或非正式：\textit{a shift in policy}）；\textbf{move}（移动——通用：\textit{move the table}）；\textbf{transfer}（调动——多指组织/官方：\textit{transfer to another department}）；\textbf{relocate}（搬迁——多指家/工作地：\textit{relocate to London}）。
\end{gram}

\begin{origin}
古英语 \textit{sciftan}（划分、安排），与古诺斯语 \textit{skipta}（划分）同源，原始日耳曼语 *\textit{skip-}。同源词 \textbf{shift}（改变；轮班；移动）、\textbf{shifty}（狡猾的；不可靠的——\textit{a shifty character}）、德语 \textit{schichten}（分层）。
\end{origin}

\vw{cottage}{/ˈkɒtɪdʒ/}{n.}

\sense{n.}{小屋，村舍}
\col{a country/seaside cottage}
\ex{They rented a small country cottage for the summer holiday.}
\col{cottage hospital 小型医院}
\ex{The village has a small cottage hospital serving the local community.}

\begin{origin}
14 世纪中古英语 \textit{cotage}，\textit{cot}（小屋）+ \textit{-age} 后缀。cot 来自古英语 \textit{cot}（小屋；棚），原始日耳曼语 *\textit{kot-}。同源词 \textbf{cot}（小屋；吊床；婴儿床）、\textbf{cottage}、\textbf{cottage industry}、德语 \textit{Kot}（茅屋）、冰岛语 \textit{kot}（小屋）。cottage 字面"小屋"——cote（鸽舍；羊栏）同源。
\end{origin}

\vp{cottage industry}

\sense{短语}{家庭手工业}
\ex{The region is known for its cottage industries, particularly hand-loomed textiles and pottery.}
\col{from cottage industry to factory}
\ex{The Industrial Revolution transformed a cottage industry into a factory system within a few decades.}

\begin{gram}{cottage industry / heavy industry / light industry 的"工业"细分}
\textbf{cottage industry}（家庭手工业——小规模、手工）；\textbf{heavy industry}（重工业——钢铁、机械）；\textbf{light industry}（轻工业——纺织、食品）；\textbf{high-tech industry}（高科技产业——信息、生物）。cottage industry 是工业革命前的典型生产形态。
\end{gram}

\end{document}
